\documentclass[12pt,a4paper]{report}
\usepackage[dvips]{graphicx}           
\usepackage[portuges]{babel}           
\usepackage{color}            
\usepackage[latin1]{inputenc}           
\usepackage[normalem]{ulem}
\setlength{\textwidth}{160mm}
\setlength{\textheight}{220mm}
\setlength{\topmargin}{-8pt}
\setlength{\oddsidemargin}{0mm}
\usepackage{amsmath,amsfonts, amssymb}
\usepackage{enumitem}% http://ctan.org/pkg/enumitem

\begin{document}

  \begin{minipage}[b]{0.1 \linewidth}
     \includegraphics[width=.9cm]{ufs_vertical_positiva.png}
  \end{minipage}\hfill	
  \begin{minipage}[b]{0.85 \linewidth}
	       \textsc{\large Universidade Federal de Sergipe}\\
	       \textsc{\large Departamento de Estat�stica e Ci�ncias Atuariais}\\
	       Disciplina: ESTAT0078 -- Infer�ncia I\\
	       Professor: Sadraque E. F. Lucena
\end{minipage}

\vspace{0.8cm}
\begin{center} {\bf Lista de Exerc�cios 12} \end{center}
\vspace{0.8cm}

	\begin{enumerate}[label=12.\arabic*)]
		
		% Quest�o 1
		\item Verifique se as distribui��es a seguir pertencem � fam�lia exponencial. Em caso afirmativo, identifique cada uma das fun��es $c(\theta)$, $T(x)$, $d(\theta)$ e $S(x)$. Para as distribui��es que pertencem � fam�lia exponencial, indique as estat�sticas suficientes.
		\begin{enumerate}
			\item $X\sim$ Poisson($\theta$), $\theta>0$.
			\item $X\sim$ Exponencial($\theta$), de forma que $E(X)=\theta$.
			\item $X\sim$ Geom�trica($\theta$), $\theta\in(0,1)$.
			\item $X\sim$ Gama($r,\theta$), $r,\theta>0$, $E(X)=r\theta$.
			\item $X\sim$ Beta($\theta,1$), $\theta>0$.
			\item $X\sim$ Normal($0,\sigma^2$), $\sigma^2>0$.
			\item $X\sim$ Normal Inversa($\mu,\lambda$), $\mu,\lambda>0$.
			\item $X\sim$ Beta($\alpha,\beta$), $\alpha,\beta>0$.
			\item $X\sim$ Uniforme($0,\theta$), $\theta>0$.
			\item $f(x) = \frac{e^{-\theta} \theta^x}{x!(1-e^{-\theta})}$, $x=1,2,\ldots,\theta>0$.
			\item $f(x) = \frac{1}{\lambda}e^{-\frac{1}{\lambda}(x-\mu)}$, $\mu\leq x<\infty$.
			\item $f(x) = \frac{\log(\theta)\, \theta^x}{\theta-1}$, $\theta>1$, $0\leq x\leq 1$.
		\end{enumerate}
		
	\end{enumerate}
	
	\newpage
	
	\noindent Respostas:
	
	\begin{enumerate}[label=12.\arabic*)]
		
		% Quest�o 1
		\item 
		\begin{enumerate}
			% a
			\item $c(\theta) = \log(\theta)$\\
			      $T(x) = x$\\
			      $d(\theta)=-\theta$\\
			      $S(x) = -\log(x!)$\\
			      Estat�stica suficiente: $T(\underset{\sim}{X}) = \sum\limits_{i=1}^n X_i$
			
			% b
			\item $c(\theta) = -\frac{1}{\theta}$\\
			      $T(x) = x$\\
			      $d(\theta)= -\log(\theta)$\\
			      $S(x) = 0$\\
			      Estat�stica suficiente: $T(\underset{\sim}{X}) = \sum\limits_{i=1}^n X_i$
			
			% c
			\item $c(\theta) = \log(1-\theta)$\\
			      $T(x) = x$\\
			      $d(\theta)= \log(\theta)$\\
			      $S(x) = 0$\\
			      Estat�stica suficiente: $T(\underset{\sim}{X}) = \sum\limits_{i=1}^n X_i$
			
			% d
			\item $c(r,\theta) = \left(-\frac{1}{\theta}, r-1\right)$\\
			      $T(x) = \left(x, \log(x)\right)$\\
			      $d(r,\theta)= -r \log(\theta) - \log(\Gamma(r))$\\
			      $S(x) = (r-1)\log(x) - \log(\Gamma(r))$\\
				  Estat�stica suficiente: $T(\underset{\sim}{X}) = \left(\sum\limits_{i=1}^n X_i, \sum\limits_{i=1}^n \log(X_i)\right)$
			
			% e
			\item $c(\theta) = \theta-1$\\
			      $T(x) = \log(x)$\\
			      $d(\theta)= \log(\theta)$\\
			      $S(x) = 0$\\
				  Estat�stica suficiente: $T(\underset{\sim}{X}) = \sum\limits_{i=1}^n \log(X_i)$
			
			% f
			\item $c(\sigma^2) = -\frac{1}{2\sigma^2}$\\
			      $T(x) = x^2$\\
			      $d(\sigma^2)= -\frac{1}{2}\log(\sigma^2)$\\
			      $S(x) = -\frac{1}{2}\log(2\pi)$\\
				  Estat�stica suficiente: $T(\underset{\sim}{X}) = \sum\limits_{i=1}^n X_i^2$
			
			% g
			\item $c(\mu,\lambda) = \left(-\frac{\lambda}{2\mu^2}, -\frac{\lambda}{2}\right)$\\
			      $T(x) = (x, \frac{1}{x})$\\
			      $d(\mu,\lambda)= \frac{\lambda}{\mu} + \frac{1}{2}\log(\lambda)$\\
			      $S(x) = -\frac{1}{2}\log(2\pi x^3)$\\
				  Estat�stica suficiente: $T(\underset{\sim}{X}) = \left(\sum\limits_{i=1}^n X_i, \sum\limits_{i=1}^n\frac{1}{X_i}\right)$
			
			% h
			\item $c(\alpha,\beta) = (\alpha-1, \beta-1)$\\
			      $T(x) = (\log(x), \log(1-x))$,\\
			      $d(\alpha,\beta)= \log(\Gamma(\alpha+\beta))-\log(\Gamma(\alpha)) - \log(\Gamma(\beta))$\\
			      $S(x) = 0$\\
				  Estat�stica suficiente: $T(\underset{\sim}{X}) = \left(\sum\limits_{i=1}^n \log(X_i), \sum\limits_{i=1}^n \log(1-X_i)\right)$
			
			% i
			\item N�o pertence.
			
			% j
			\item $c(\theta) = \log(\theta)$\\
			      $T(x) = x$\\
			      $d(\theta)= -\theta-\log(1-e^{-\theta})$\\
			      $S(x) = -\log(x!)$\\
				  Estat�stica suficiente: $T(\underset{\sim}{X}) = \sum\limits_{i=1}^n X_i$
			
			% k
			\item N�o pertence
			
			% l
			\item $c(\theta) = \log(\theta)$\\
			      $T(x) = x$\\
			      $d(\theta)= \log(\log(\theta))-\log(\theta-1)$\\
			      $S(x) = 0$\\
				  Estat�stica suficiente: $T(\underset{\sim}{X}) = \sum\limits_{i=1}^n X_i$
		\end{enumerate}
		
	\end{enumerate}
	
\end{document}